\section{Output Analysis}
\label{sec:output}

This is a \textbf{terminating} simulation: the system starts empty and idle and
ends when the 50\textsuperscript{th} batch ships, so no warm-up period is deleted.
Each scenario is run as $n$ independent replications with a unique random seed per
replication, using a \texttt{Parameter Variation} experiment that writes one row of
KPIs per replication to CSV.

\paragraph{Number of replications.}
We size $n$ with the half-width method. For a KPI with sample mean $\bar x$ and
sample standard deviation $s$ over $n_0$ replications, the number of replications
needed for a half-width $\varepsilon$ at confidence $1-\alpha$ is
\begin{equation}
  n \;\ge\; \left(\frac{t_{\,n_0-1,\,1-\alpha/2}\,s}{\varepsilon}\right)^{\!2},
  \label{eq:halfwidth}
\end{equation}
with $\alpha=0.05$ and a target relative precision $\varepsilon \le \SI{5}{\percent}$
of the mean. We size on \textbf{mean flow time}, the noisiest of the reported
operational KPIs; weekly throughput has lower relative variability and is therefore
not the binding KPI for sizing.

We ran $n=20$ replications of the baseline. Table~\ref{tab:reps} reports the achieved
half-widths and the $n$ that Eq.~\eqref{eq:halfwidth} requires for each KPI.

\begin{table}[H]
\centering
\caption{Replication sizing from the $n=20$ baseline. Required $n$ targets a
\SI{5}{\percent} relative half-width at \SI{95}{\percent} confidence.}
\label{tab:reps}
\begin{tabularx}{\textwidth}{X r r r r r}
\toprule
KPI & mean $\bar x$ & st.\ dev $s$ & half-width & rel.\ h-w & required $n$ \\
\midrule
Mean flow time (h)         & 180.0  & 7.30   & 3.42   & \SI{1.9}{\percent} & 3 \\
Completion / makespan (h)  & 436.2  & 14.35  & 6.71   & \SI{1.5}{\percent} & 2 \\
Throughput (bottles)       & 60{,}972 & 2{,}281 & 1{,}067 & \SI{1.8}{\percent} & 3 \\
Avg.\ WIP (batches)        & 19.5   & 0.71   & 0.33   & \SI{1.7}{\percent} & 3 \\
Dryer utilisation          & 0.949  & 0.005  & 0.002  & \SI{0.2}{\percent} & 1 \\
Scrap rate                 & 0.062  & 0.032  & 0.015  & \SI{24}{\percent}  & 6\textsuperscript{*} \\
\bottomrule
\end{tabularx}
\end{table}

\noindent\textsuperscript{*}Scrap rate has a small mean, so a \SI{5}{\percent}
\emph{relative} target is not meaningful; we instead size it on an \emph{absolute}
half-width of $\pm 0.03$, which requires $n=6$. The achieved absolute half-width at
$n=20$ is $\pm 0.015$ (i.e.\ $\pm 1.5$ percentage points).

\paragraph{Chosen replication count.}
Every reported KPI needs at most six replications to reach the target precision, so
$n=20$ leaves a comfortable margin and is used for all scenarios. All means in
Section~\ref{sec:experiments} are reported with their \SI{95}{\percent} confidence
intervals on this basis.

\paragraph{Resolving differences between scenarios.}
For scenario comparisons the relevant precision is the confidence interval on the
\emph{difference} from the baseline, not on a single mean. At $n=20$ this resolves
the large effects (the Industry 4.0 scrap reduction, and the combined
dryer-plus-cleaning-team configuration on flow time and makespan), while the effect
of any single added machine on throughput falls below the detectable threshold. As
shown in Section~\ref{sec:experiments}, that non-detectability is itself a result:
shipped volume is fixed by the 50-batch plan and per-batch yield, so adding one unit
of capacity shifts the bottleneck and changes flow time and utilisation, but does not
move throughput by a statistically meaningful amount. The only lever that raises
shipped output is the scrap reduction, which acts on yield rather than capacity.
